%% APPENDIX E -- radio-frequency interference screening.
%% Owner: appendix-triage.  Carries \label{app:rfi}.
\section{Radio-frequency interference screening}
\label{app:rfi}

Interference is the first alternative explanation for an unattributed
crossing, so the screening this survey has and the coverage it lacks are both
set out here. Millimetre interferometry is an unusually poor
environment for terrestrial interference, though not an empty one: ITU Radio
Regulations Article~5 \citep{ITU2024} carries active allocations throughout
90--275\,GHz and WRC-19 identified parts of 275--450\,GHz for land-mobile and
fixed use. The specific hazard is the \RadarLo--\RadarHi\,GHz Earth
exploration-satellite allocation used by spaceborne cloud radar
\citep{CloudSat2002,EarthCARE2023}, a documented contaminant of ALMA Band~3;
\RadarNWin{} of the \RadarNTot{} census windows overlaps it, carrying
\RadarNCross{} threshold crossings and \RadarNStage{} rank-passing events, so
no result here rests on it. What favours this band is that terrestrial
transmitters are sparse compared with centimetre wavelengths and, decisively,
that an interferometer is not a total-power detector: a signal not in the far
field at the phase centre acquires a rapidly varying geometric phase across
the array, and is suppressed in the correlation rather than imaged.

\emph{The coverage of the screening is uneven.} The allocation test above and
the sky- and intermediate-frequency tests below are computed for every
searched window of the census, as is the \NCtrl-position rank screen, which
asks whether an excess is at the star or everywhere in the field. Recurrence
reaches only the crossings with a covering repeat block, and the visibility
fit is a consistency check rather than a filter. The correlator flagging is
the observatory's own, replayed from each block's delivered calibration
script, and was not retained window by window. \emph{No per-channel
interference flag of our own exists, and none can be recovered.} Four
measurements follow, and each could have failed.

\emph{First, clustering in sky frequency.} If interference dominated the
crossing population, the crossings would pile up at particular sky
frequencies. They do cluster --- but on the survey's own molecular lines.
The test needs a measured peak frequency in the frame an interfering carrier
would stand still in, so it runs topocentrically over the \AppMNCrossF{} of
the \EvNCrossRaw{} crossings for which one was measured, against the
transitions the search was tuned to. The other \CsNCrossNoFreq{} are carbon
monoxide toward \CsCrossNoFreqStar, \MxBpNoFreqClause{} of its
\MxBpNCoWord{} such crossings; for those the catalogue carries only the
velocity offset from the line, so the frequency Table~\ref{tab:ledger} prints
is recovered from that offset rather than measured, and a frequency
reconstructed from a line offset cannot be used to ask whether frequencies
crowd onto lines.
Of the \AppMNCrossF, \AppMTubeObs{} fall within
$\pm\DispoHalfKms$\,km\,s$^{-1}$ of a searched transition against
\AppMTubeNull{} expected if each window's peak fell anywhere on that
window's own channel grid, $\pv{\AppMTubeP}$. The densest band,
\AppMBandLo--\AppMBandHi\,GHz, holds \AppMBandN{} of them, \AppMBandNAttr{}
within that velocity window of CO($2{\to}1$); and once the line-attributed
crossings are set aside, the \AppMUnattrN{} that remain show no clustering at
all ($\pv{\AppMUnattrP}$). \emph{The concentration is the attribution, not a
residue left over after it.} Table~\ref{tab:ledger} counts the same band as
\EvTubeLedgerN{} crossings with \EvTubeLedgerAttr{} attributed, over a
different set --- all \EvNCrossRaw{} of them, each attributed in its own
star's rest frame --- and the \MxExtraWord{} rows it adds are the ones this
test cannot hold: \MxExtraClause.

\emph{Second, the same test against the atmosphere.} A residual of the
atmospheric calibration stands still in topocentric frequency exactly as a
transmitter does, and it is narrow where the atmosphere has lines.
\TlNLine{} transitions of ozone, water and oxygen lie within
\TlFreqLo--\TlFreqHi\,GHz \citep{Pickett1998}, a median \TlSpacing\,GHz
apart. No crossing of either sample falls within \TlOffMin\,MHz of one; the
median separation is \TlOffMed\,GHz; and \TlTubeObs{} crossing lies inside the
same $\pm\TlTubeHalf$\,km\,s$^{-1}$ window used above, against
\TlTubeNull{} expected from the windows' own channel grids. The one
concentration a reader will notice --- \TlBandAll{} crossings of the census
and the hold-out between \TlBandLo{} and \TlBandHi\,GHz --- is not atmospheric
either, because that interval holds \TlBandNLine{} telluric transitions;
\TlBandObs{} of the \TlBandNSub{} unattributed crossings whose window the
catalogue tabulates fall in it against \TlBandNull$\,\pm\,$\TlBandNullSd{}
expected from where the spectral windows were tuned: the carbon monoxide
tuning, seen from the other side.

\emph{Third, fixed intermediate frequency.} An artefact of the signal chain
would sit at a fixed position within the spectral window rather than at a
fixed sky frequency. This test is run over every window for which the
catalogue releases a peak frequency --- \IfPopNWin{} of the \NWindows{}, of
which only \IfPopNCross{} carry a threshold crossing --- because such an artefact would
occupy the same intermediate frequency in all of them and not only where the
trigger fires. It does not: the peaks' fractional positions within their own
windows have median \RfiFracMed{} and are uniform.

\emph{Fourth, a cross-target occupancy screen over every window the pipeline
extracted, not only the crossings.} That set is wider than the census:
\RfiOccNWin{} windows toward \RfiOccNStars{} stars, being the census
together with the rows its audit removed, the reserved hold-out, the archival
extension, the out-of-sample control, and \OqOtherWin{} windows in
\OqOtherEb{} blocks that entered none of them. Interference belongs to the
observatory and not to the star, so a terrestrial carrier should recur at the
same sky frequency toward \emph{unrelated} stars. It does not, and the
sharpest place to look is the \MxNRankWord{} unattributed crossings that
outrank every control position in their own window, \MxRankStars, both of
which carry a measured frequency. Taking them in that order,
\MxRankFreqs\,GHz are covered by \MxOtherWinRange{} further windows, toward
\MxOtherStarsRange{} other stars in \MxOtherBlocksRange{} other blocks, and
\MxSameChanClause. With each window's peak randomised
on its own channel grid, the number of cross-star coincidences within
\RfiOccTol\,MHz is \RfiOccObs{} against
\RfiOccNull$\,\pm\,$\RfiOccNullSd{} expected ($\pv{\RfiOccP}$); among the
\RfiOccSigN{} windows whose peak exceeds the trigger it is \RfiOccSigObs{}
against \RfiOccSigNull$\,\pm\,$\RfiOccSigNullSd, i.e.\ \emph{below} the
null, and with the molecular-line mask applied \RfiOccMskObs{} against
\RfiOccMskNull$\,\pm\,$\RfiOccMskNullSd. There is no excess to explain. The
mask has to be applied first, because \RfiOccMaskN{} of those \RfiOccSigN{}
windows sit on a masked transition (\RfiOccSpeciesList), and a carbon
monoxide line falls at nearly the same sky frequency toward every nearby
star --- exactly the coincidence the screen is looking for.

\emph{One execution block fails these tests, and it fails them as a block
rather than as a carrier.} Its \DqBlockNCross{} crossings toward
\LgVerRejectStar{} are the highest statistics in the survey outside
$\beta$~Pictoris and HD~48370, and \S\ref{sec:bdblock} identifies the
failure. Whether interference accounts for it is what belongs here, and three
readings say not: the visibility fit returns a large imaginary part, placing
the emission away from the phase centre; the block's \BdNSibling{} sibling
executions of the same field, array and spectral setup produce no crossing at
all; and each of the \DqBlockNCross{} frequencies is covered by
\DqOtherWinLo--\DqOtherWinHi{} further windows toward
\DqOtherStarLo--\DqOtherStarHi{} other stars, of which \DqOtherSame{} carries
a crossing in the same channel.

\emph{One reading does point at a mechanism, and it can be tested across the
survey.} \DqBlockIfMult{} of the \DqBlockNCross{} sit at the identical
channel of two different spectral windows, which is the signature of a fixed
intermediate frequency; and a feature with no geometric fringe phase --- a
spur common to every antenna --- does not fringe-track, so it gathers at the
phase centre rather than at the star. Here the two are a
twentieth of an arcsecond apart, so the block cannot test its own hypothesis;
the survey can, because the star lies \EvPcOffLo--\EvPcOffHi\,arcsec from the
phase centre in \EvPcNWin{} windows whose control annulus brackets it,
putting a median of \EvPcNNear{} probes within two synthesised beams of the
phase centre and the star well outside. Against probes at the same distance
from the star but elsewhere in azimuth --- the radius has to be matched, or a
centrally concentrated disc would imitate the effect --- those probes are
higher by $\EvPcDiff\pm\EvPcErr$ in the statistic, consistent with zero and
bounded by \EvPcBound, smaller by a factor of about \EvPcFactor{} than the
\EvPcNeed{} by which this block's ring is raised. \emph{The phase centre is
not where the excess lives.} The evidence supports a data-quality failure of
one execution block, and the block criterion of
Appendix~\ref{sec:blockquality} --- which fails \BqNFail{} block of the
\BqNBlocks{} searched --- is what disposes of these \EvNCrossFlag{} crossings
(\S\ref{sec:bdblock}).

\emph{None of this substitutes for the ON/OFF cadence a dedicated search
uses.} The claim is only that \emph{no
identified interference mechanism accounts for the unattributed events}:
allocation tables cannot exclude harmonics, intermodulation products,
unconventional emitters or near-field sources, and archival data without a
simultaneous off-source strategy cannot do better. The four tests establish
that the crossing population carries none of the signatures interference
leaves --- a statement about the population, not a clearance of any one
crossing.

%% "The null result rests on internal validation alone" lives once, in the
%% discussion's lessons subsection, and not here.  fig:appm is not included
%% -- it compared our own two statistics and the paper states only the one it
%% uses -- so figures/appm_power.pdf and the \AppMPowerOld* and
%% \RfiOcc{Prior,Dir,Stale}* families will be retired.
%%
%% ★ r16-mask: the cross-target screen of the fourth test is now scoped to
%% the ADOPTED attribution (round 660, maskrfi_v660.py), which makes it two
%% events and not one.  round 49's \RfiEventLead, \RfiFreqThem,
%% \RfiSameChanClause, \RfiNotTestableClause, \RfiOtherWin,
%% \RfiOtherStarsRange and \RfiOtherBlocksRange are therefore unreferenced
%% and will be retired; rfi_v399.py still reproduces them and its grammar
%% selftest still runs, so nothing breaks if the scope is ever restored.
